\section{Investment Implications}

\subsection{Break-Even Analysis}

Table~\ref{tab:break-even} presents break-even analysis for the five highest-redundancy-value
corridors, computing the capital cost of the best practical alternate, the annual benefit
(redundancy value), and the net present value at a 7\% discount rate over 30 years.

\subsubsection{Donner Pass: I-70W Completion}

A new trans-Sierra interstate crossing (hypothetical I-70W, following a route roughly
parallel to US-50 through El Dorado County) would require approximately 120 miles of new
construction through mountainous terrain. Based on recent mountain interstate construction
costs (\$175M/mile for the new Coalfields Expressway precedent; \$130M/mile for recent
Colorado I-70 reconstruction), an estimate of \$21 billion is applied.

\begin{itemize}
\item Capital cost: \$21B
\item Annual benefit (redundancy value): \$1.9B/yr
\item Simple payback: 11.1 years
\item NPV at 7\% over 30 years: $\text{PV}(\$1.9\text{B}/\text{yr}, 30\text{ yr}, 7\%) -
\$21\text{B} = \$1.9\text{B} \times 12.41 - \$21\text{B} = \$23.6\text{B} - \$21\text{B}
= +\$2.6\text{B}$
\end{itemize}

The I-70W NPV (\$+2.6B at 7\%) is positive but modest. Sensitivity: at 5\% discount rate,
NPV = \$+5.7B; at 10\%, NPV = $-\$2.1B$ (negative). The I-70W investment is rate-sensitive:
it clears the federal discount rate (3.5--7\% range for USDOT project evaluation) but
does not have large NPV cushion at the high end of that range.

\textbf{Donner snowshed program (hardening without redundancy)}:
The Donner Summit snowshed option --- covering the critical 14-mile Donner Summit section
with a permanent reinforced concrete shelter structure, eliminating avalanche and extreme
snowfall closures --- is estimated at \$800M (based on the Connaught Creek Tunnel precedent
in British Columbia, scaled to 14-mile coverage and US cost factors). Annual benefit:
reduces closure frequency from 50/yr to approximately 30/yr; annual cost reduction =
$\$2.4\text{B} \times (20/50) = \$960\text{M/yr}$. Simple payback: 0.83 years. NPV:
$\$960\text{M} \times 12.41 - \$800\text{M} = +\$11.1\text{B}$. \textbf{This is a
substantially better short-term investment than I-70W}: lower capital cost, faster payback,
higher NPV at the 7\% rate. The snowshed is the correct first investment at Donner.

\subsubsection{Dallas Interchange: US-287 Upgrade}

Upgrading US-287 (Fort Worth to Wichita Falls, approximately 140 miles) to full interstate
standard (4 lanes, controlled access, interstate geometry), reducing the Dallas alternate
detour from 200 to 80 miles:

\begin{itemize}
\item Capital cost: \$3.0B (140 miles $\times$ \$21.4M/mile for rural interstate upgrade,
including controlled access acquisition)
\item Annual benefit (redundancy value): \$0.4B/yr
\item Simple payback: 7.5 years
\item NPV at 7\%: $\$0.4\text{B} \times 12.41 - \$3.0\text{B} = +\$1.97\text{B}$
\end{itemize}

\subsubsection{I-95 Baltimore: Managed Lane Program (Cross-Reference to E.1)}

The I-95 Baltimore--Washington corridor has a different investment profile. The primary
cost driver is not closure events (B1 is already moderate at 1.20 multiplier) but a
combination of recurring closure costs (\$0.6B/yr) and ordinary congestion delay (the
Baltimore and Washington beltways are ATRI top-10 bottlenecks). The managed lane program
analyzed in E.1 \citep{ROUTE_D1} addresses both:

\begin{itemize}
\item Capital cost of managed lanes (E.1 estimate): \$12B
\item Annual benefit from closure cost reduction: \$0.4B/yr
\item Annual benefit from congestion delay reduction: \$2.0B/yr (from ATRI bottleneck data)
\item Total annual benefit: \$2.4B/yr
\item Simple payback: 5.0 years
\item NPV at 7\%: $\$2.4\text{B} \times 12.41 - \$12\text{B} = +\$17.8\text{B}$
\end{itemize}

I-95 Baltimore has the highest NPV in the ROUTE corpus precisely because it benefits from
\textit{both} closure cost reduction \textit{and} congestion relief --- a compound
investment benefit.

\subsection{Current IIJA Funding Alignment}

The Infrastructure Investment and Jobs Act (IIJA, 2021) \citep{IIJA2021} established two
funding streams relevant to the corridors analyzed here:

\textbf{PROTECT Program} (\$8.7B, 5 years): Climate resilience for highways. Current
eligibility is based on total SFHA miles and bridge vulnerability ratings. As argued in
D.1, the consecutive-miles metric should be added to PROTECT eligibility. The closure
cost data from D.2 provides the economic justification for PROTECT funding scale: if the
top-5 corridors impose \$5.4B/year in closure costs, the \$8.7B PROTECT authorization
represents less than two years of current closure costs --- and does not include
redundancy investments, only hardening.

\textbf{National Highway Performance Program} (NHPP, \$36B, 5 years): General highway
investment, including congestion relief. NHPP does not have a redundancy value criterion.
Adding redundancy value as an investment criterion would redirect a portion of NHPP
funding toward alternate route investments (US-287 upgrade, Stevens Pass improvement)
that score poorly on V/C-based bottleneck criteria but have strong redundancy-value NPV.

\begin{table}[ht]
\centering
\caption{Break-even analysis for top-5 redundancy-value investments. Capital cost,
annual benefit (redundancy value or combined closure+congestion benefit), simple
payback, and NPV at 7\% over 30 years. I-95 Baltimore includes managed lane
congestion relief benefit from E.1.}
\label{tab:break-even}
\begin{tabular}{@{}lrrccc@{}}
\toprule
Corridor & Capital & Annual & Payback & NPV & NPV \\
         & Cost    & Benefit & (yr) & (7\%, 30yr) & (5\%, 30yr) \\
         & (\$B)   & (\$B/yr) &    & (\$B) & (\$B) \\
\midrule
Donner snowshed          & 0.8  & 0.96 &  0.8 & +11.1 & +13.0 \\
Dallas US-287 upgrade    & 3.0  & 0.40 &  7.5 &  +1.97 & +3.3 \\
I-95 Baltimore (mgd lanes) & 12.0 & 2.40 & 5.0 & +17.8 & +25.6 \\
Donner I-70W alternate   & 21.0 & 1.90 & 11.1 &  +2.6 & +5.7 \\
Snoqualmie US-2 upgrade  &  2.1 & 0.23 &  9.1 &  +0.75 & +1.4 \\
\bottomrule
\end{tabular}
\end{table}

\subsection{Optimal Investment Sequencing}

The compound investment principle established in Section~\ref{sec:redundancy} and the
break-even data in Table~\ref{tab:break-even} support the following investment sequence
for Track D corridors:

\begin{enumerate}
\item \textbf{Donner snowshed}: NPV \$+11.1B at 0.83-year payback. Build first. This
eliminates the highest-frequency closure events (the 30+ moderate storms per year) while
the larger alternate-route investment (I-70W) is authorized and designed.
\item \textbf{I-95 Baltimore managed lanes}: NPV \$+17.8B, addresses both closure costs
and congestion. Second priority on NPV grounds; requires E.1 managed lane program.
\item \textbf{Gulf Coast I-10 hardening}: D1 = 8.4 (2024) rising to 9.1 (2050).
Hardening cost estimated \$18--25B for full-corridor elevation (out of scope for D.2;
see E.2 for I2.0 framework design). Must be authorized by 2038 to be operational before
2050 exposure threshold.
\item \textbf{Dallas US-287 upgrade}: NPV \$+1.97B; solid investment case; lower priority
than items 1--3 on NPV per dollar.
\item \textbf{I-70W Donner alternate}: NPV \$+2.6B; correct eventual investment but
capital-intensive and rate-sensitive; sequence after snowshed delivers near-term benefit.
\end{enumerate}
